\section{Discussion and Future Scope}
\label{sec:discussion}

The results of our analysis and the development of the PrimeFusion framework indicate that blockchain technology, when thoughtfully adapted, holds significant promise for enhancing the security and resilience of UAV communication networks. The successful integration of a lightweight consensus mechanism (PoC) and efficient data compression (CBOR) demonstrates that it is possible to overcome the most prohibitive barriers—latency and resource overhead—that have traditionally hindered the application of blockchain in this domain. The PoC model, in particular, represents a paradigm shift away from computationally expensive consensus towards a model that is synergistically aligned with the cooperative nature of multi-UAV missions. By incentivizing cooperation, the framework not only secures the network but also reinforces the collaborative behaviors essential for mission success.

However, the limitations discussed in Section \ref{sec:limitations} highlight that blockchain is not a panacea. The trade-offs between security, latency, and computational cost must be carefully managed. The applicability of a framework like PrimeFusion will be context-dependent, with the greatest benefits realized in scenarios where decentralized trust and data integrity are paramount, such as in multi-agency disaster response operations or military applications where a central point of failure is unacceptable. For less critical applications or those operating under extreme resource constraints, the overhead of a blockchain may not be justifiable.

Building upon the foundation of the PrimeFusion framework, several promising avenues for future research emerge:
\begin{itemize}
    \item \textbf{Machine Learning Integration for Adaptive Consensus:} Future work could explore the integration of machine learning, particularly reinforcement learning, to dynamically tune the parameters of the PoC consensus mechanism. An AI-driven consensus could adapt to changing network conditions, mission priorities, and threat landscapes in real-time, further optimizing the balance between security and performance.
    \item \textbf{Hierarchical Blockchain Architectures and Sharding:} To address the scalability limitations for ultra-dense swarms, hierarchical or sharded blockchain architectures could be investigated. In such a model, UAVs could be organized into clusters, each with its own local blockchain, with a higher-level blockchain used to record inter-cluster transactions. This would distribute the consensus load and reduce the communication overhead on individual nodes.
    \item \textbf{Integration with 5G/6G and Edge Computing:} The synergy between blockchain, next-generation wireless networks (5G/6G), and mobile edge computing (MEC) is a rich area for exploration. 5G/6G can provide the high-bandwidth, low-latency communication required for blockchain operations, while MEC servers can offload the most computationally intensive tasks, such as smart contract execution or cryptographic operations, from the UAVs themselves.
    \item \textbf{Real-World Experimental Validation:} While the simulation results presented for the PrimeFusion framework are promising, the ultimate validation of its effectiveness requires deployment and testing in a real-world field environment. Such experiments would provide invaluable data on the framework's performance under actual flight conditions, channel impairments, and hardware constraints.
\end{itemize}
By pursuing these research directions, the capabilities of blockchain-enabled UAV networks can be further expanded, paving the way for more autonomous, intelligent, and secure aerial systems.
